\] To prove it, choose each orbit's representative of minimum frequency radius. There are at most \(8R^2\) oriented orbits meeting the radius-R box, each with \(|q_A|\leq C_AR^2\). Cauchy--Schwarz bounds their part of (ND.13) divided by \(4\pi^2\) by \(\sqrt8 C_AR^3\sqrt\Gamma\). For all other orbits, every frequency has radius greater than R, and the sum is bounded by \(C_AR^{-s}\mathcal H_s(h)\leq J/2\). Rearranging proves (ND.17). This bound is for a fixed regular symbol with a nonzero obstruction. It is not uniform over all mixed-sign Fourier polynomials.

\subsection{A one-sign Fourier theorem without a bandwidth restriction}\label{nondegeneracy:a-one-sign-fourier-theorem-without-a-bandwidth-restriction}

\textbf{Theorem 3.} Let \(A\in\mathrm{SL}(2,\mathbb Z)\) be hyperbolic. If the nonconstant coefficients of a real even smooth mean-zero function \(g=\sum_v c_vb_v\) all have the same sign and at least one is nonzero, then \[
                             \Gamma_A(g)>0.                 \tag{ND.18}
\] This is a statement about cosine coefficients, not pointwise nonnegativity of a mean-zero function. There is no restriction on the number of signed orbits.

\textbf{Proof.} Change the overall sign if necessary, so \(c_v\geq0\), and put \[
 M=\sum_v c_vvv^{\mathsf T}=-\operatorname{Hess}g(0)/(8\pi^2).
\] Polarization of \(q_A\) gives the exact discriminant identity \[
 B_A(u,v)^2-q_A(u)q_A(v)=\frac\Delta4\det(u,v)^2.            \tag{ND.19}
\] Expanding \(\mathcal E_A(M)\) and \(\det M\) shows \[
       \mathcal E_A(M)\geq\left(\sum_v c_v|q_A(v)|\right)^2>0.
 \tag{ND.20}
\] For clarity, each cross coefficient in the difference between the two sides is zero or positive: if \(q_A(u)q_A(v)\geq0\), its extra term is \(\Delta\det(u,v)^2\); if the signs differ, it is \(4B_A(u,v)^2\) by (ND.19). The diagonal terms agree. Also \(q_A(v)\ne0\) for a nonzero integer v: otherwise A would have a rational eigenline and a rational algebraic-unit eigenvalue, hence eigenvalue \(1\) or \(-1\), contrary to hyperbolicity. Absolute smooth convergence justifies the sums. Equation (ND.12), or (ND.14), now proves (ND.18). End of proof.

For an even finite Fourier polynomial, there is a useful exact identity behind this criterion. If \(a_{\mathcal O}\) denotes the cosine sum of h on a signed orbit, then \[
 \boxed{\sum_{\mathcal O}q_A(v_{\mathcal O})a_{\mathcal O}(h)
                         =-2\pi t\mathcal E_A(M).}          \tag{ND.21}
\] It follows from (ND.13)--(ND.14) using \(\mathcal L_A b_v(0)=-8\pi^2q_A(v)\) and \(\mathcal E_A(-8\pi^2M)=64\pi^4\mathcal E_A(M)\). Thus this test excludes a \emph{nonzero coboundary bracket}, not just an identically zero bracket.

A coarse finite-bandwidth constant is explicit. Define \[
 m_A=\gcd(|c|,|d-a|,|b|),\quad C_A=|c|+|d-a|+|b|,\quad
 D_A=1+\max\{|a|+|b|,|c|+|d|\}.
\] For the one-sign class supported in \(\|v\|_\infty\leq R\), \[
 \boxed{\displaystyle
 \Gamma_A(g)\geq
 \frac{\pi^2t^2m_A^4}{2C_A^2D_A^6}\,R^{-6}\|g\|_2^4.}      \tag{ND.22}
\] Indeed, all generated frequencies lie in the radius-\(D_AR\) box. Choosing representatives from those generated frequencies gives \(\sum_{\mathcal O}q_A(v_{\mathcal O})^2\leq4C_A^2(D_AR)^6\). Cauchy--Schwarz in (ND.21), (ND.1), (ND.20), and \(|q_A(v)|\geq m_A\), \(\|g\|_2^2=2\sum c_v^2\), prove (ND.22). For the Pell matrix this specializes to \[
                \Gamma_T(g)\geq\frac{\pi^2}{32768R^6}\|g\|_2^4.
 \tag{ND.23}
